    \chapter{Multivariable Calculus}
    \section{Vector-valued functions}
    \begin{definitionenv}
        Suppose $f_1,\cdots,f_n$ are functions fron $\R$ to $\R$. We call $$\bm F(t) = (f_1(t), f_2(t), \cdots, f_n(t))$$ a \textbf{vector-valued function} and we define the \textbf{limits,derivatives} and \textbf{integrals} by the following:
        $$ \begin{aligned}
            \lim_{t\to p} \bm F(t) &= \left( \lim_{t\to p} f_1(t), \cdots, \lim_{t\to p} f_n(t) \right),\\
            \bm F'(t) &= (f_1'(t), \cdots, f_n'(t)),\\
            \int_a^b \bm F(t) \dd t &= \left( \int_a^b f_1(t) \dd t, \cdots, \int_a^b f_n(t) \dd t \right).
        \end{aligned} $$
    \end{definitionenv}
    \begin{theoremenv}
        Suppose $\bm F$ and $\bm G$ are vector-valued functions from $\R$ to $\R^n$ and $u$ is a real-valued function from $\R$ to $\R$. If $F,G$ and $u$ are differentiable on in interval, then so are $F+G$, $uF$ and $F\cdot G$ and we have
        \begin{enumerate}[label=(\arabic*)]
            \item $(\bm F+\bm G)' = \bm F' + \bm G'$;
            \item $(u\bm F)' = u' \bm F + u \bm F'$;
            \item $(\bm F\cdot \bm G)' = \bm F'\cdot \bm G + \bm F\cdot \bm G'$.
        \end{enumerate}
        For the case $n=3$, then $\bm F\times \bm G$ is also differentiable with $(\bm F\times \bm G)' = \bm F'\times \bm G + \bm F \times \bm G'$.
    \end{theoremenv}
    \begin{propositionenv}
        Suppose $\bm F$ is a vector-valued function from $\R$ to $\R^n$. If $\bm F$ is differentiable and be constant length/norm on an open interval $I$, i.e., $\| \bm F(t)\| = c$ for any $t\in I$, then $\bm F \cdot \bm F' = 0$ on $I$. Namely, $\bm F'(t)$ is always perpendicular to $\bm F(t)$ for any $t\in I$.
    \end{propositionenv}
    \begin{proofenv}
        Let $g(t) := \bm F(t)\cdot \bm F(t) = \| \bm F(t)\|^2 = c^2$ for any $t\in I$. Then $g'(t) = 2 \bm F(t)\cdot \bm F'(t) = 0$ for any $t\in I$. Thus $\bm F(t)\cdot \bm F'(t) = 0$ for any $t\in I$.
    \end{proofenv}
    \begin{theoremenv}
        Suppose $\bm F:\R \to \R^n$ and $u:\R \to \R$. Let $\bm G:=\bm F \circ u$. Then:
        \begin{enumerate}[label=(\arabic*)]
            \item If $u$ is continuous at $t$ and $\bm F$ is continuous at $u(t)$, then $\bm G$ is continuous at $t$.
            \item If the derivative $u'(t)$ and $\bm F'(u(t))$ exist, then $\bm G'(t)$ exists and $\bm G'(t) = u'(t) \bm F'(u(t))$.
        \end{enumerate}
    \end{theoremenv}
    \begin{theoremenv}
        Suppose $\bm F:\R \to \R^n$ is continuous on $[a,b]$, then:
        \begin{enumerate}[label=(\arabic*)]
            \item $\bm F$ is integrable on $[a,b]$;
            \item For any $c\in [a,b]$, we can define the primitive or indefinite integral function $\bm A(x) := \displaystyle \int_c^x \bm F(t) \dd t$ for any $x\in [a,b]$. Then $\bm A$ is differentiable on $(a,b)$ and $\bm A'(x) = \bm F(x)$ for any $x\in (a,b)$.
            \item If the derivative $\bm F'(t)$ is continuous on some open interval $I$, then for every $c$ and $x$ in $I$, we have
            $$ \bm F(x) = \bm F(c) + \int_c^x \bm F'(t) \dd t. $$
        \end{enumerate}
    \end{theoremenv}
    \begin{theoremenv}
        Suppose $\bm F$ and $\bm G$ are from $\R$ to $\R^n$ and are integrable on $[a,b]$. Then :
        \begin{enumerate}[label=(\arabic*)]
            \item for any $\alpha , \beta \in \R$, $\alpha \bm F + \beta \bm G$ is also integrable on $[a,b]$ with
            $$ \int_a^b (\alpha \bm F + \beta \bm G)(t) \dd t = \alpha \int_a^b \bm F(t) \dd t + \beta \int_a^b \bm G(t) \dd t; $$
            \item Let $C \in \R^n$, then the dot product $C \cdot \bm F$ is also integrable on $[a,b]$ with
            $$ \int_a^b (C\cdot \bm F)(t) \dd t = C \cdot \int_a^b \bm F(t) \dd t. $$
            \item If $\| \bm F \|$ is also integrable on $[a,b]$, then 
            $$ \left\| \int_a^b \bm F(t) \dd t \right\| \le \int_a^b \| \bm F(t)\| \dd t. $$
        \end{enumerate}
    \end{theoremenv}
    \begin{proofenv}
        We only prove (3) here. Let $\bm A:= \displaystyle \int_a^b \bm F(t) \dd t$. The case where $\left\|\bm A\right\| = 0 $ is trivial. Suppose $\left\|\bm A\right\| > 0$. Then
        $$ \|\bm A\|^2 = \bm A \cdot \bm A = \int_a^b \bm F(t) \cdot \bm A \dd t \le \int_a^b \| \bm F(t)\| \| \bm A\| \dd t = \| \bm A\| \int_a^b \| \bm F(t)\| \dd t. $$
        Thus $\displaystyle \left\|\int_a^b \bm F(t) \dd t\right\| = \| \bm A\| \le \int_a^b \| \bm F(t)\| \dd t$.
    \end{proofenv}
    \section{Curves}
    \begin{definitionenv}
        Suppose $\bm X$ is a vector-valued function from $\R$ to $\R^n$ with $n=2$ or $n=3$. Suppose the domain of $\bm X$ is an interval $D\subseteq \R$. If $\bm X$ is continuous on $D$, then we call the range/image of $\bm X$ a \textbf{curve} in $\R^n$. We also say that the curve is \textbf{described} (parametrically) by the function $\bm X$ with parameter $t\in D$. We also call $D$ a \textbf{parameter interval}.
    \end{definitionenv}
    \begin{remarkenv}
        The curve described by $\bm X$ is invarient under the change of parameter. For example, if $u:\R \to \R$ and $I = \{u(t) :t \in J\}$ for some interval $J$, then $\bm Y := \bm X\circ u$ also describes the same curve as $\bm X$. We say $\bm X$ and $\bm Y$ are \textbf{equivalent} in describing the curve $C$.
    \end{remarkenv}
    \begin{definitionenv}
        Let $C$ be a curve described by $\bm X$ (note $\bm X$ is continuous). If the derivative $\bm X'(t)$ exists and $\bm X'(t) \ne 0$, then the line through point $\bm X(t)$ with direction vector $\bm X'(t)$ is called the \textbf{tangent line} of the curve $C$ at point $\bm X(t)$.
    \end{definitionenv}
    \begin{remarkenv}
        Though the derivative of $\bm X$ and $\bm Y=\bm X\circ u$ are different, the tangent lines of the curve $C$ at point $\bm X(t) = \bm Y(s)$ are the same since $\bm Y'(s) = u'(s) \bm X'(u(s))$ is parallel to $\bm X'(u(s))$. Thus the tangent line of a curve at a point is well-defined and does not depend on the choice of parameter.
    \end{remarkenv}
    \begin{definitionenv}
        Consider a motion (of a particle) described by a vector-valued function $\bm X: \R \to \R^n$ with $n=2$ or $n=3$. The derivative $\bm X'(t)$ is called the \textbf{velocity} at time $t$. $\| \bm X'(t)\|$ is called the \textbf{speed} at time $t$. The second derivative $\bm X''(t)$ is called the \textbf{acceleration} at time $t$. In some textbooks, we use $\bm{r}(t)$ for $\bm X(t)$, $\bm{v}(t)$ for $\bm X'(t)$, $v(t)$ for $\| \bm X'(t)\|$ and $\bm{a}(t)$ for $\bm X''(t)$.
    \end{definitionenv}
    \begin{exampleenv}
        Suppose $\bm{r}(t) = (r \cos(f(t)), r \sin(f(t)), bf(t))$ with $r>0$ and $b\ne 0$ fixed. This is called a \textbf{helix/helical curve} in $\R^3$. Note a helix is not contained in any plane in $\R^3$ but it is contained fully in a cylinder/cylindrical surface in $\R^3$.
    \end{exampleenv}
    \begin{definitionenv}
        Consider a curve described by $\bm X(t)$. We define the \textbf{unit tangent vector} $\bm T$ as $$\bm T(t) := \dfrac{\bm X'(t)}{\| \bm X'(t)\|} \quad \text{whenever } \bm X'(t) \ne 0. $$ Note that $\| \bm T(t) \| = 1$ and thus $\bm T(t) \cdot \bm T'(t) = 0$. We define the \textbf{principal normal vetor} $\bm N$ as
        $$ \bm N(t) := \frac{\bm T'(t)}{\| \bm T'(t)\|} \quad \text{whenever } \bm T'(t) \ne 0. $$
        Note $\left\|\bm N(t)\right\| =1$ and $\bm N(t) \cdot \bm T(t) = 0$. Furthermore, we call the plane through the point $\bm X(t)$ and spanned by $\bm T(t)$ and $\bm N(t)$ the \textbf{osculating plane} of the curve at point $\bm X(t)$. 
    \end{definitionenv}
    \begin{remarkenv}
        Consider three points $\bm X(t_1), \bm X(t_2)$ and $\bm X(t_3)$ on the curve. As $t_2 $ and $t_3$ approach $t_1$, the plane determined by $\bm X(t_1), \bm X(t_2)$ and $\bm X(t_3)$, supposing the three points are not collinear, approaches the osculating plane at point $\bm X(t_1)$. Note if the curve is a planar curve, then the osculating plane is the fixed plane containing the curve.
    \end{remarkenv}
    \begin{theoremenv}
        For a motion described by $\bm r(t)$, the acceleration $\bm a(t) := \bm r''(t)$ is a linear combination of $\bm T(t)$ and $\bm T'(t)$ given by $$ \bm a(t) = v'(t) \bm T(t) + v(t) \bm T'(t). $$
        If, additionally $\bm T'(t) \ne 0$, then
        $$ \bm a(t) = v'(t) \bm T(t) + v(t) \| \bm T'(t)\| \bm N(t). $$
    \end{theoremenv}
    \begin{remarkenv}
        We call $v'(t)$ and $v(t) \| \bm T'(t)\|$ the tangential and normal components of the acceleration $\bm a(t)$, respectively. In some textbook, the \textbf{binormal vector} $\bm B$ is also introduced as $B(t) := \bm T(t) \times \bm N(t)$. Note $\left\| \bm B(t) \right\|=1$ and $\bm T(t), \bm N(t)$ and $\bm B(t)$ together provide a moving reference frame for $\R^3$ along the curve.
    \end{remarkenv}
    \begin{definitionenv}
        For a curve in $\R^2$ and $\R^3$ described by $\bm r(t)$ for $t\in [a,b]$ (Note $\bm r$ is required to be continuous on $[a,b]$). Let $P=\{t_0, t_1, \ldots, t_n\}$ be a partition of $[a,b]$. We denote the polygon whose vertices are $\bm r(t_0), \bm r(t_1), \ldots, \bm r(t_n)$ by $\Pi (P)$ and denote the \textbf{length of polygon} by $\| \Pi (P)\|$. Then
        $$ \| \Pi (P)\| = \sum_{i=1}^n \| \bm r(t_i) - \bm r(t_{i-1})\|. $$
        If there exists an $M>0$ such that $\| \Pi (P)\| \le M$ for any partition $P$ of $[a,b]$, then we say the curve is said to be \textbf{rectifiable} and its arc length, denoted by $\Lambda(a,b)$, is defined as
        $$ \Lambda(a,b) := \sup \{ \| \Pi (P)\| : P \text{ is a partition of } [a,b]\}. $$
        If there's no such $M$, then the curve is said to be \textbf{non-rectifiable}.
    \end{definitionenv}
    \begin{theoremenv}
        For a curve in $\R^2$ or $\R^3$ described by $\bm r(t)$ for $t\in [a,b]$. If $\bm v(t) := \bm r'(t)$ is continuous on $[a,b]$, then the curve is rectifiable and its arc length satisfies
        $$ \Lambda(a,b) \le \int_a^b \| \bm v(t)\| \dd t. $$

    \end{theoremenv}
    \begin{proofenv}
        Suppose $P = \{t_0, t_1, \ldots, t_n\}$ is a partition of $[a,b]$. Then
        $$ \| \Pi (P)\| = \sum_{i=1}^n \| \bm r(t_i) - \bm r(t_{i-1})\| = \sum_{i=1}^n \left\| \int_{t_{i-1}}^{t_i} \bm v(t) \dd t \right\| \le \sum_{i=1}^n \int_{t_{i-1}}^{t_i} \| \bm v(t)\| \dd t = \int_a^b \| \bm v(t)\| \dd t. $$
        Since $\Lambda (a,b) $ is the least upper bound, we have $\Lambda(a,b)\displaystyle \le \int_a^b \| \bm v(t)\| \dd t$.
    \end{proofenv}
    \begin{theoremenv}
        For a rectifiable curve in $\R^2$ or $\R^3$ described by $\bm r(t)$ for $t\in [a,b]$. If $a<c<b$ and $C_1$ and $C_2$ are the curves corresponding to $\bm r(t)$ on $[a,c]$ and $[c,b]$, respectively, then $C_1$ and $C_2$ are both rectifiable and $\Lambda(a,b) = \Lambda(a,c) + \Lambda(c,b)$.
    \end{theoremenv}
    \begin{proofenv}
        Suppose $P_1$ and $P_2$ are partitions of $[a,c]$ and $[c,b]$, respectively. Let $P = P_1 \cup P_2$ be the joint partition. Then $ \| \Pi (P_1)\| + \| \Pi (P_2)\| = \| \Pi (P)\| \le \Lambda(a,b)$. Thus $\| \Pi (P_1)\|$ and $\| \Pi (P_2)\|$ are both bounded above by $\Lambda(a,b)$ for any partitions $P_1$ and $P_2$. Thus $C_1$ and $C_2$ are both rectifiable.
        \begin{itemize}
            \item For a fixed $P_2$, $\| \Pi (P_1)\|  \le \Lambda(a,b) - \| \Pi (P_2)\|$ holds for any partition $P_1$ of $[a,c]$ and there by $$\Lambda(a,c) \le \Lambda(a,b) - \| \Pi (P_2)\| \implies \| \Pi (P_2) \| \le \Lambda(a,b) - \Lambda(a,c).$$
            \item since the above inequality holds for any partition $P_2$ of $[c,b]$, we have $\Lambda(a,b) \ge \Lambda(a,c) + \Lambda(c,b)$.
        \end{itemize}
        We still need to show $\Lambda(a,b) \le \Lambda(a,c) + \Lambda(c,b)$. Suppose $P$ is a partition of $[a,b]$. Adding in $c$ into $P$, we obtain a partition $P_1$ of $[a,c]$ and a partition $P_2$ of $[c,b]$ such that
        $$ \| \Pi (P)\| \le \| \Pi (P_1)\| + \| \Pi (P_2)\| \le \Lambda(a,c) + \Lambda(c,b). $$
        Thus $\Lambda(a,b) \le \Lambda(a,c) + \Lambda(c,b)$. Combining the two inequalities, we have $\Lambda(a,b) = \Lambda(a,c) + \Lambda(c,b)$.
    \end{proofenv}
    \begin{definitionenv}
        Suppose a curve in $\R^2$ or $\R^3$ is described by $\bm r(t)$ for $t\in [a,b]$. We define the \textbf{arc-length function} $s$ as
        $$ s(t) := \begin{cases}
            0, & t=a; \\
            \Lambda(a,t), & t\in (a,b].
        \end{cases}
        $$
    \end{definitionenv}
    \begin{theoremenv}
        For any rectifiable curve, the arc-length function $s$ is monotonically increasing on $[a,b]$.
    \end{theoremenv}
    \begin{theoremenv}
        Suppose a rectifiable curve described by $\bm r(t)$ for $t\in [a,b]$. Let $s$ denote the arc-length function. If the speed $v(t) = \| \bm r'(t)\|$ is continuous on $[a,b]$, then the derivative $s'(t)$ exists and is given by $s'(t) = v(t)$ for any $t\in [a,b]$.
    \end{theoremenv}
    \begin{proofenv}
        Suppose $h>0$, then
        $$ \frac{s(t+h) - s(t)}{h} \le \frac1h \int_t^{t+h} v(u) \dd u \to v(t) \quad \text{as } h\to 0^+. $$
        We also have
        $$ \frac{s(t+h) - s(t)}{h} \ge \left\|\frac{\bm r(t+h) - \bm r(t)}{h}\right\| \to v(t) \quad \text{as } h\to 0^+. $$
        By the squeeze theorem, we know $\displaystyle \lim_{h\to 0^+} \frac{s(t+h) - s(t)}{h} = v(t)$ and thus $s'(t) = v(t)$. The case where $h<0$ can be proved similarly.
    \end{proofenv}
    \begin{remarkenv}
        Note
        $$ s(t_2) - s(t_1) = \int_{t_1}^{t_2} v(t) \dd t,$$
        provided that $v(t)$ is continuous on $[a,b]$. Particularly, $$\Lambda(a,b) = s(b) - s(a) = \int_a^b v(t) \dd t.$$
    \end{remarkenv}
    \begin{exampleenv}
        Consider the curve in $\R^2$ given by $\bm r(t) = (t, f(t))$ for $t\in [a,b]$ where $f$ is a scalar-valued function differentiable on $[a,b]$. Then the arc-length function is given by
        $$ s(t) = \int_a^t v(u) \dd u = \int_a^t \sqrt{1 + (f'(u))^2} \dd u. $$
    \end{exampleenv}
    \begin{definitionenv}
        Let $\bm T$ denote the unit tangent vector and $s$ denote the arc-length function of a rectifiable curve described by $\bm r(t)$ for $t\in [a,b]$. We call $\dfrac{\dd \bm T}{\dd s}$ the \textbf{curvature vector} of the curve provided it's well-defined. Note
        $$ \frac{\dd \bm T}{\dd s} = \frac{\dd \bm T}{\dd t} \cdot \frac{\dd t}{\dd s} = \frac{\bm T'(t)}{v(t)} = \frac{\left\|\bm T'(t)\right\| \bm N(t)}{v(t)}. $$
        This means that the curvature vector is parallel to the principal normal vector $\bm N$. We call $$\kappa(t) := \left\| \dfrac{\dd \bm T}{\dd s}\right\| = \dfrac{\left\|\bm T'(t)\right\|}{v(t)}$$ the \textbf{curvature number} of the curve at point $\bm r(t)$.
    \end{definitionenv}
    \begin{definitionenv}
        When $\kappa(t) \ne 0$, we call $\rho(t) := \dfrac{1}{\kappa(t)}$ the \textbf{radius of curvature}. Then then circle in the osculating plane with radius $\rho(t)$ and centered at $\bm r(t) + \dfrac{1}{\kappa(t)} \bm N(t)$ is called the \textbf{osculating circle}. Note the osculating circle is tangent to the curve at point $\bm r(t)$ and has the same curvature as the curve at point $\bm r(t)$.
    \end{definitionenv}

    \begin{theoremenv}
        For motion along a curve in $\R^2$ or $\R^3$ with velocity $\bm v(t)$, speed $v(t)$, acceleration $\bm a(t)$ and curvature $\kappa$. Then we have
        $$ \bm a(t) = v'(t) \bm T(t) + \kappa (t) v^2(t) \bm N(t). $$
        and
        $$ \kappa (t) = \frac{\left\|\bm a (t) \times \bm v (t)\right\|}{v^3 (t)}
        $$
    \end{theoremenv}
    \section{Curves in Polar Coordinates}
    Sometimes it's more natural to use polar coordinates to describe a curve in $\R^2$:
    $$ \bm r = (r \cos \theta, r \sin \theta) = r (\cos \theta, \sin \theta). $$
    \begin{definitionenv}
        We define the unit radial vector $\bm e_r$ and the tangential vector $\bm e_\theta$ as
        $$ \bm e_r = (\cos \theta, \sin \theta) \quad \text{and} \quad \bm e_\theta = \dv{\bm e_r}{\theta} = (-\sin \theta, \cos \theta). $$
        Note $\| \bm e_r\| = \| \bm e_\theta\| = 1$ and $\bm e_r \cdot \bm e_\theta = 0$.
    \end{definitionenv}
    \begin{exampleenv}
        Let $\bm r(t) = f(t)$, $\theta(t) = g(t)$ for $t\in [a,b]$. Then
        $$ \bm r(t) = r \bm e_r = f(t) \bm e_r  + 0 \cdot \bm e_\theta. $$
        Thus the velocity $\bm v(t)$ is given by
        $$ \bm v(t) = \dv{\bm r}{t} = \dv{r}{t} \bm e_r + r \dv{\bm e_r}{t} = \dv{r}{t} \bm e_r + r \dv{\bm e_r}{\theta} \dv{\theta}{t} = f'(t) \bm e_r + \left[ f(t) g'(t)  \right]\bm e_\theta. $$
        We call $f'(t)$ the \textbf{radial component} of the velocity $\bm v(t)$ and $f(t) g'(t)$ the \textbf{transverse component} of the velocity $\bm v(t)$.
        Note the speed $v(t)$ is given by
        $$v(t) =\sqrt{ \bm v(t) \cdot \bm v(t)} = \sqrt{(f'(t))^2 + (f(t) g'(t))^2} $$
        and the acceleration $\bm a(t)$ is given by
        $$\bm a(t) = \left[ \dv[2]{r}{t} - r\left( \dv{\theta}{t} \right)^2 \right] \bm e_r + \left[ r \dv[2]{\theta}{t} + 2 \dv{r}{t} \dv{\theta}{t} \right] \bm e_\theta = \left[ \dv[2]{r}{t} - r\left( \dv{\theta}{t} \right)^2 \right] \bm e_r + \frac1r \frac{\dd}{\dd t} \left( r^2 \dv{\theta}{t} \right) \bm e_\theta .$$
        We say the acceleration $\bm a(t)$ is \textbf{radical} if the transverse component is always zero, i.e., $\displaystyle r^2 \dv{\theta}{t}$ is a constant. 

    \end{exampleenv}

    \newpage